\section{Internal Hom}
\label{section-internal-hom}

\noindent
Let $\mathcal{C}$ be a category and let $\mathcal{O}$ be a presheaf
of rings. Let $\mathcal{F}$, $\mathcal{G}$ be presheaves of
$\mathcal{O}$-modules. Consider the rule
$$
U \longmapsto \Hom_{\mathcal{O}_U}(\mathcal{F}|_U, \mathcal{G}|_U).
$$
For $\varphi : V \to U$ in $\mathcal{C}$ we define a restriction mapping
$$
\Hom_{\mathcal{O}_U}(\mathcal{F}|_U, \mathcal{G}|_U)
\longrightarrow
\Hom_{\mathcal{O}_V}(\mathcal{F}|_V, \mathcal{G}|_V)
$$
by restricting via the relocalization morphism
$j : \mathcal{C}/V \to \mathcal{C}/U$, see
Sites, Lemma \ref{sites-lemma-relocalize}. Hence this defines a
presheaf $\SheafHom_\mathcal{O}(\mathcal{F}, \mathcal{G})$.
In addition, given an element
$\varphi \in \Hom_{\mathcal{O}|_U}(\mathcal{F}|_U, \mathcal{G}|_U)$
and a section $f \in \mathcal{O}(U)$ then we can define
$f\varphi \in \Hom_{\mathcal{O}|_U}(\mathcal{F}|_U, \mathcal{G}|_U)$
by either precomposing with multiplication by $f$ on $\mathcal{F}|_U$
or postcomposing with multiplication by $f$ on $\mathcal{G}|_U$ (it gives
the same result). Hence we in fact get a presheaf of $\mathcal{O}$-modules.
There is a canonical ``evaluation'' morphism
$$
\mathcal{F}
\otimes_{p, \mathcal{O}}
\SheafHom_\mathcal{O}(\mathcal{F}, \mathcal{G})
\longrightarrow
\mathcal{G}.
$$

\begin{lemma}
\label{lemma-internal-hom}
If $\mathcal{C}$ is a site, $\mathcal{O}$ is a sheaf of rings,
$\mathcal{F}$ is a presheaf of $\mathcal{O}$-modules, and
$\mathcal{G}$ is a sheaf of $\mathcal{O}$-modules, then
$\SheafHom_\mathcal{O}(\mathcal{F}, \mathcal{G})$
is a sheaf of $\mathcal{O}$-modules.
\end{lemma}

\begin{proof}
Omitted. Hints: Note first that
$\SheafHom_\mathcal{O}(\mathcal{F}, \mathcal{G})
= \SheafHom_\mathcal{O}(\mathcal{F}^\#, \mathcal{G})$, which reduces
the question to the case where both $\mathcal{F}$ and $\mathcal{G}$
are sheaves. The result for sheaves of sets is
Sites, Lemma \ref{sites-lemma-glue-maps}.
\end{proof}

\begin{lemma}
\label{lemma-internal-hom-restriction}
Let $(\mathcal{C}, \mathcal{O})$ be a ringed site.
Let $\mathcal{F}, \mathcal{G}$ be sheaves of $\mathcal{O}$-modules.
Then formation of $\SheafHom_\mathcal{O}(\mathcal{F}, \mathcal{G})$
commutes with restriction to $U$ for $U \in \Ob(\mathcal{C})$.
\end{lemma}

\begin{proof}
Immediate from the definition.
\end{proof}

\begin{remark}
\label{remark-pullback-internal-hom}
Let $f : (\mathcal{C}, \mathcal{O}_\mathcal{C}) \to
(\mathcal{D}, \mathcal{O}_\mathcal{D})$ be a morphism
of ringed sites. Let
$\mathcal{F}, \mathcal{G}$ be sheaves of $\mathcal{O}_\mathcal{D}$-modules.
There is a canonical map
$$
f^*\SheafHom_{\mathcal{O}_\mathcal{D}}(\mathcal{F}, \mathcal{G})
\longrightarrow
\SheafHom_{\mathcal{O}_\mathcal{C}}(f^*\mathcal{F}, f^*\mathcal{G})
$$
Namely, this map is adjoint to the map
$$
\SheafHom_{\mathcal{O}_\mathcal{D}}(\mathcal{F}, \mathcal{G})
\longrightarrow
f_*\SheafHom_{\mathcal{O}_\mathcal{C}}(f^*\mathcal{F}, f^*\mathcal{G})
$$
defined as follows. Say $f$ is given by the continuous functor
$u : \mathcal{D} \to \mathcal{C}$. For sections over $V \in \Ob(\mathcal{D})$
we use the map
\begin{align*}
\Gamma(V, \SheafHom_{\mathcal{O}_\mathcal{D}}(\mathcal{F}, \mathcal{G}))
& =
\Hom_{\mathcal{O}_V}(\mathcal{F}|_V, \mathcal{G}|_V) \\
& \longrightarrow
\Hom_{\mathcal{O}_{u(V)}}(f^*\mathcal{F}|_{u(V)}, f^*\mathcal{G}|_{u(V)}) \\
& =
\Gamma(u(V), \SheafHom_{\mathcal{O}_\mathcal{C}}(f^*\mathcal{F},
f^*\mathcal{G})) \\
& =
\Gamma(V, f_*\SheafHom_{\mathcal{O}_\mathcal{C}}(f^*\mathcal{F},
f^*\mathcal{G}))
\end{align*}
where for the arrow we use pullback by the morphism
$(\mathcal{C}/u(V), \mathcal{O}_{u(V)}) \to
(\mathcal{D}/V, \mathcal{O}_V)$ induced by $f$.
\end{remark}

\noindent
In the situation of Lemma \ref{lemma-internal-hom}
the ``evaluation'' morphism
factors through the tensor product of sheaves of modules
$$
\mathcal{F}
\otimes_\mathcal{O}
\SheafHom_\mathcal{O}(\mathcal{F}, \mathcal{G})
\longrightarrow
\mathcal{G}.
$$

\begin{lemma}
\label{lemma-internal-hom-commute-limits}
Internal hom and (co)limits.
Let $\mathcal{C}$ be a category and let $\mathcal{O}$ be a presheaf of rings.
\begin{enumerate}
\item For any presheaf of $\mathcal{O}$-modules $\mathcal{F}$ the functor
$$
\textit{PMod}(\mathcal{O}) \longrightarrow \textit{PMod}(\mathcal{O})
, \quad
\mathcal{G} \longmapsto \SheafHom_\mathcal{O}(\mathcal{F}, \mathcal{G})
$$
commutes with arbitrary limits.
\item For any presheaf of $\mathcal{O}$-modules $\mathcal{G}$ the functor
$$
\textit{PMod}(\mathcal{O}) \longrightarrow \textit{PMod}(\mathcal{O})^{opp}
, \quad
\mathcal{F} \longmapsto \SheafHom_\mathcal{O}(\mathcal{F}, \mathcal{G})
$$
commutes with arbitrary colimits, in a formula
$$
\SheafHom_\mathcal{O}(\colim_i \mathcal{F}_i, \mathcal{G})
=
\lim_i \SheafHom_\mathcal{O}(\mathcal{F}_i, \mathcal{G}).
$$
\end{enumerate}
Suppose that $\mathcal{C}$ is a site, and $\mathcal{O}$ is a sheaf of rings.
\begin{enumerate}
\item[(3)] For any sheaf of $\mathcal{O}$-modules $\mathcal{F}$ the functor
$$
\textit{Mod}(\mathcal{O}) \longrightarrow \textit{Mod}(\mathcal{O})
, \quad
\mathcal{G} \longmapsto \SheafHom_\mathcal{O}(\mathcal{F}, \mathcal{G})
$$
commutes with arbitrary limits.
\item[(4)] For any sheaf of $\mathcal{O}$-modules $\mathcal{G}$ the functor
$$
\textit{Mod}(\mathcal{O}) \longrightarrow \textit{Mod}(\mathcal{O})^{opp}
, \quad
\mathcal{F} \longmapsto \SheafHom_\mathcal{O}(\mathcal{F}, \mathcal{G})
$$
commutes with arbitrary colimits, in a formula
$$
\SheafHom_\mathcal{O}(\colim_i \mathcal{F}_i, \mathcal{G})
=
\lim_i \SheafHom_\mathcal{O}(\mathcal{F}_i, \mathcal{G}).
$$
\end{enumerate}
\end{lemma}

\begin{proof}
Let $\mathcal{I} \to \textit{PMod}(\mathcal{O})$, $i \mapsto \mathcal{G}_i$
be a diagram. Let $U$ be an object of the category $\mathcal{C}$.
As $j_U^*$ is both a left and a right adjoint we see that
$\lim_i j_U^*\mathcal{G}_i = j_U^* \lim_i \mathcal{G}_i$.
Hence we have
\begin{align*}
\SheafHom_\mathcal{O}(\mathcal{F}, \lim_i \mathcal{G}_i)(U)
& =
\Hom_{\mathcal{O}_U}(\mathcal{F}|_U, \lim_i \mathcal{G}_i|_U) \\
& =
\lim_i \Hom_{\mathcal{O}_U}(\mathcal{F}|_U, \mathcal{G}_i|_U) \\
& = \lim_i \SheafHom_\mathcal{O}(\mathcal{F}, \mathcal{G}_i)(U)
\end{align*}
by definition of a limit. This proves (1). Part (2) is proved in exactly the
same way. Part (3) follows from (1) because the limit of a diagram of sheaves
is the same as the limit in the category of presheaves.
Finally, (4) follow because, in the formula we have
$$
\Mor_{\textit{Mod}(\mathcal{O})}(
\colim_i \mathcal{F}_i, \mathcal{G})
=
\Mor_{\textit{PMod}(\mathcal{O})}(
\colim^{PSh}_i \mathcal{F}_i, \mathcal{G})
$$
as the colimit $\colim_i \mathcal{F}_i$ is the sheafification of
the colimit $\colim^{PSh}_i \mathcal{F}_i$ in
$\textit{PMod}(\mathcal{O})$. Hence (4) follows from (2)
(by the remark on limits above again).
\end{proof}

\begin{lemma}
\label{lemma-internal-hom-exact}
Let $(\mathcal{C}, \mathcal{O})$ be a ringed site.
Let $\mathcal{F}$, $\mathcal{G}$ be $\mathcal{O}$-modules.
\begin{enumerate}
\item If $\mathcal{F}_2 \to \mathcal{F}_1 \to \mathcal{F} \to 0$
is an exact sequence of $\mathcal{O}$-modules, then
$$
0 \to
\SheafHom_\mathcal{O}(\mathcal{F}, \mathcal{G}) \to
\SheafHom_\mathcal{O}(\mathcal{F}_1, \mathcal{G}) \to
\SheafHom_\mathcal{O}(\mathcal{F}_2, \mathcal{G})
$$
is exact.
\item If $0 \to \mathcal{G} \to \mathcal{G}_1 \to \mathcal{G}_2$
is an exact sequence of $\mathcal{O}$-modules, then
$$
0 \to
\SheafHom_\mathcal{O}(\mathcal{F}, \mathcal{G}) \to
\SheafHom_\mathcal{O}(\mathcal{F}, \mathcal{G}_1) \to
\SheafHom_\mathcal{O}(\mathcal{F}, \mathcal{G}_2)
$$
is exact.
\end{enumerate}
\end{lemma}

\begin{proof}
Follows from Lemma \ref{lemma-internal-hom-commute-limits} and
Homology, Lemma \ref{homology-lemma-exact-functor}.
\end{proof}

\begin{lemma}
\label{lemma-internal-hom-adjoint-tensor}
Let $\mathcal{C}$ be a category. Let $\mathcal{O}$ be a presheaf of
rings.
\begin{enumerate}
\item Let $\mathcal{F}$, $\mathcal{G}$, $\mathcal{H}$ be
presheaves of $\mathcal{O}$-modules. There is a canonical isomorphism
$$
\SheafHom_\mathcal{O}
(\mathcal{F} \otimes_{p, \mathcal{O}} \mathcal{G}, \mathcal{H})
\longrightarrow
\SheafHom_\mathcal{O}
(\mathcal{F}, \SheafHom_\mathcal{O}(\mathcal{G}, \mathcal{H}))
$$
which is functorial in all three entries (sheaf Hom in
all three spots). In particular,
$$
\Mor_{\textit{PMod}(\mathcal{O})}(
\mathcal{F} \otimes_{p, \mathcal{O}} \mathcal{G}, \mathcal{H})
=
\Mor_{\textit{PMod}(\mathcal{O})}(
\mathcal{F}, \SheafHom_\mathcal{O}(\mathcal{G}, \mathcal{H}))
$$
\item
Suppose that $\mathcal{C}$ is a site, $\mathcal{O}$ is a sheaf of rings,
and $\mathcal{F}$, $\mathcal{G}$, $\mathcal{H}$ are sheaves of
$\mathcal{O}$-modules. There is a canonical isomorphism
$$
\SheafHom_\mathcal{O}
(\mathcal{F} \otimes_\mathcal{O} \mathcal{G}, \mathcal{H})
\longrightarrow
\SheafHom_\mathcal{O}
(\mathcal{F}, \SheafHom_\mathcal{O}(\mathcal{G}, \mathcal{H}))
$$
which is functorial in all three entries (sheaf Hom in
all three spots). In particular,
$$
\Mor_{\textit{Mod}(\mathcal{O})}(
\mathcal{F} \otimes_\mathcal{O} \mathcal{G}, \mathcal{H})
=
\Mor_{\textit{Mod}(\mathcal{O})}(
\mathcal{F}, \SheafHom_\mathcal{O}(\mathcal{G}, \mathcal{H}))
$$
\end{enumerate}
\end{lemma}

\begin{proof}
This is the analogue of
Algebra, Lemma \ref{algebra-lemma-hom-from-tensor-product}.
The proof is the same, and is omitted.
\end{proof}

\begin{lemma}
\label{lemma-tensor-commute-colimits}
Tensor product and colimits.
Let $\mathcal{C}$ be a category and let $\mathcal{O}$ be a presheaf of rings.
\begin{enumerate}
\item For any presheaf of $\mathcal{O}$-modules $\mathcal{F}$ the functor
$$
\textit{PMod}(\mathcal{O}) \longrightarrow \textit{PMod}(\mathcal{O})
, \quad
\mathcal{G} \longmapsto \mathcal{F} \otimes_{p, \mathcal{O}} \mathcal{G}
$$
commutes with arbitrary colimits.
\item
Suppose that $\mathcal{C}$ is a site, and $\mathcal{O}$ is a sheaf of rings.
For any sheaf of $\mathcal{O}$-modules $\mathcal{F}$ the functor
$$
\textit{Mod}(\mathcal{O}) \longrightarrow \textit{Mod}(\mathcal{O})
, \quad
\mathcal{G} \longmapsto \mathcal{F} \otimes_\mathcal{O} \mathcal{G}
$$
commutes with arbitrary colimits.
\end{enumerate}
\end{lemma}

\begin{proof}
This is because tensor product is adjoint to internal hom according
to Lemma \ref{lemma-internal-hom-adjoint-tensor}.
See Categories, Lemma \ref{categories-lemma-adjoint-exact}.
\end{proof}

\begin{lemma}
\label{lemma-adjoint-hom-restrict}
Let $\mathcal{C}$ be a category, resp.\ a site
Let $\mathcal{O} \to \mathcal{O}'$ be a map of presheaves, resp.\ sheaves
of rings. Then
$$
\Hom_\mathcal{O}(\mathcal{G}, \mathcal{F}) =
\Hom_{\mathcal{O}'}(\mathcal{G},
\SheafHom_\mathcal{O}(\mathcal{O}', \mathcal{F}))
$$
for any $\mathcal{O}'$-module $\mathcal{G}$ and $\mathcal{O}$-module
$\mathcal{F}$.
\end{lemma}

\begin{proof}
This is the analogue of
Algebra, Lemma \ref{algebra-lemma-adjoint-hom-restrict}.
The proof is the same, and is omitted.
\end{proof}

\begin{lemma}
\label{lemma-j-shriek-and-tensor}
Let $(\mathcal{C}, \mathcal{O})$ be a ringed site.
Let $U \in \Ob(\mathcal{C})$.
For $\mathcal{G}$ in $\textit{Mod}(\mathcal{O}_U)$
and $\mathcal{F}$ in $\textit{Mod}(\mathcal{O})$
we have $j_{U!}\mathcal{G} \otimes_\mathcal{O} \mathcal{F} =
j_{U!}(\mathcal{G} \otimes_{\mathcal{O}_U} \mathcal{F}|_U)$.
\end{lemma}

\begin{proof}
Let $\mathcal{H}$ be an object of $\textit{Mod}(\mathcal{O})$.
Then
\begin{align*}
\Hom_\mathcal{O}(j_{U!}(\mathcal{G} \otimes_{\mathcal{O}_U} \mathcal{F}|_U),
\mathcal{H})
& =
\Hom_{\mathcal{O}_U}(\mathcal{G} \otimes_{\mathcal{O}_U} \mathcal{F}|_U,
\mathcal{H}|_U) \\
& =
\Hom_{\mathcal{O}_U}(\mathcal{G},
\SheafHom_{\mathcal{O}_U}(\mathcal{F}|_U, \mathcal{H}|_U)) \\
& =
\Hom_{\mathcal{O}_U}(\mathcal{G},
\SheafHom_\mathcal{O}(\mathcal{F}, \mathcal{H})|_U) \\
& =
\Hom_\mathcal{O}(j_{U!}\mathcal{G},
\SheafHom_\mathcal{O}(\mathcal{F}, \mathcal{H})) \\
& =
\Hom_\mathcal{O}(j_{U!}\mathcal{G} \otimes_\mathcal{O} \mathcal{F},
\mathcal{H})
\end{align*}
The first equality because $j_{U!}$ is a left adjoint to restriction
of modules.
The second by Lemma \ref{lemma-internal-hom-adjoint-tensor}.
The third by Lemma \ref{lemma-internal-hom-restriction}.
The fourth because $j_{U!}$ is a left adjoint to restriction
of modules.
The fifth by Lemma \ref{lemma-internal-hom-adjoint-tensor}.
The lemma follows from this and the Yoneda lemma.
\end{proof}

\begin{remark}
\label{remark-j-shriek-tensor}
Let $\mathcal{C}$ be a site. Let $\mathcal{F}$ be a sheaf of
sets on $\mathcal{C}$ and consider the localization morphism
$j : \Sh(\mathcal{C})/\mathcal{F} \to \Sh(\mathcal{C})$.
See Sites, Definition \ref{sites-definition-localize-topos}.
We claim that (a) $j_!\mathbf{Z} = \mathbf{Z}_\mathcal{F}^\#$
and (b) $j_!(j^{-1}\mathcal{H}) = j_!\mathbf{Z} \otimes_\mathbf{Z} \mathcal{H}$
for any abelian sheaf $\mathcal{H}$ on $\mathcal{C}$.
Let $\mathcal{G}$ be an abelian on $\mathcal{C}$.
Part (a) follows from the Yoneda lemma because
$$
\Hom(j_!\mathbf{Z}, \mathcal{G}) =
\Hom(\mathbf{Z}, j^{-1}\mathcal{G}) =
\Hom(\mathbf{Z}_\mathcal{F}^\#, \mathcal{G})
$$
where the second equality holds because both sides of
the equality evaluate to the set of maps from $\mathcal{F} \to \mathcal{G}$
viewed as an abelian group. For (b) we use the Yoneda lemma and
\begin{align*}
\Hom(j_!(j^{-1}\mathcal{H}), \mathcal{G})
& =
\Hom(j^{-1}\mathcal{H}, j^{-1}\mathcal{G}) \\
& =
\Hom(\mathbf{Z}, \SheafHom(j^{-1}\mathcal{H}, j^{-1}\mathcal{G})) \\
& =
\Hom(\mathbf{Z}, j^{-1}\SheafHom(\mathcal{H}, \mathcal{G})) \\
& =
\Hom(j_!\mathbf{Z}, \SheafHom(\mathcal{H}, \mathcal{G})) \\
& =
\Hom(j_!\mathbf{Z} \otimes_\mathbf{Z} \mathcal{H}, \mathcal{G})
\end{align*}
Here we use adjunction, the fact that taking $\SheafHom$ commutes
with localization, and Lemma \ref{lemma-internal-hom-adjoint-tensor}.
\end{remark}

\begin{lemma}
\label{lemma-hom-finite-presentation-colimit}
Let $(\mathcal{C}, \mathcal{O})$ be a ringed site.
Let $\mathcal{F}$ be an $\mathcal{O}$-module of finite presentation. Let
$\mathcal{G} = \colim_{\lambda \in \Lambda} \mathcal{G}_\lambda$
be a filtered colimit of $\mathcal{O}$-modules. Then the canonical map
$$
\colim_\lambda \SheafHom_\mathcal{O}(\mathcal{F}, \mathcal{G}_\lambda)
\longrightarrow
\SheafHom_\mathcal{O}(\mathcal{F}, \mathcal{G})
$$
is an isomorphism.
\end{lemma}

\begin{proof}
It suffices to show the arrow is an isomorphism after restriction to $U$
for all $U$ in $\mathcal{C}$. Both taking colimits of sheaves of
modules and taking internal hom commute with restriction to $U$.
See for example Lemmas \ref{lemma-exactness-pushforward-pullback} and
\ref{lemma-internal-hom-restriction}. Fix $U$. Given a covering
$\{U_i \to U\}_{i \in I}$, then it suffices to prove the restriction
to each $U_i$ is an isomorphism. Hence we may assume
$\mathcal{F}$ has a global presentation
$$
\bigoplus\nolimits_{j = 1, \ldots, m}
\mathcal{O}
\longrightarrow
\bigoplus\nolimits_{i = 1, \ldots, n}
\mathcal{O}
\to
\mathcal{F}
\to
0
$$
The functor $\SheafHom_\mathcal{O}(-, -)$ commutes with finite direct sums
in either variable and $\SheafHom_\mathcal{O}(\mathcal{O}, -)$
is the identity functor. By this and by Lemma \ref{lemma-internal-hom-exact}
we obtain an exact sequence
$$
0 \to
\SheafHom_\mathcal{O}(\mathcal{F}, \mathcal{G}) \to
\bigoplus\nolimits_{i = 1, \ldots, n} \mathcal{G} \to
\bigoplus\nolimits_{j = 1, \ldots, m} \mathcal{G}
$$
Since filtered colimits are exact in $\textit{Mod}(\mathcal{O})$
by Lemma \ref{lemma-limits-colimits}
also the top row in the following commutative diagram is exact
$$
\xymatrix{
0 \ar[r] &
\colim_\lambda \SheafHom_\mathcal{O}(\mathcal{F}, \mathcal{G}_\lambda)
\ar[r] \ar[d] &
\colim_\lambda \bigoplus\nolimits_{i = 1, \ldots, n} \mathcal{G}_\lambda
\ar[r] \ar[d] &
\colim_\lambda \bigoplus\nolimits_{j = 1, \ldots, m} \mathcal{G}_\lambda
\ar[d] \\
0 \ar[r] &
\SheafHom_\mathcal{O}(\mathcal{F}, \mathcal{G}) \ar[r] &
\bigoplus\nolimits_{i = 1, \ldots, n} \mathcal{G} \ar[r] &
\bigoplus\nolimits_{j = 1, \ldots, m} \mathcal{G}
}
$$
Since the right two vertical arrows are isomorphisms we conclude.
\end{proof}

\begin{lemma}
\label{lemma-finite-presentation-quasi-compact-colimit}
Let $(\mathcal{C}, \mathcal{O})$ be a ringed site.
Let $\mathcal{G} = \colim_{\lambda \in \Lambda} \mathcal{G}_\lambda$
be a filtered colimit of $\mathcal{O}$-modules.
Let $\mathcal{F}$ be an $\mathcal{O}$-module of finite presentation.
Then we have
$$
\colim_\lambda \Hom_\mathcal{O}(\mathcal{F}, \mathcal{G}_\lambda)
=
\Hom_\mathcal{O}(\mathcal{F}, \mathcal{G}).
$$
if the hypotheses of
Sites, Lemma \ref{sites-lemma-directed-colimits-global-sections} part (4)
are satisfied for the site $\mathcal{C}$; please see
Sites, Remark \ref{sites-remark-stronger-conditions}.
\end{lemma}

\begin{proof}
Set $\mathcal{H} =
\SheafHom_\mathcal{O}(\mathcal{F}, \colim \mathcal{G}_\lambda)$
and $\mathcal{H}_\lambda =
\SheafHom_\mathcal{O}(\mathcal{F}, \mathcal{G}_\lambda)$.
Recall that
$$
\Hom_\mathcal{O}(\mathcal{F}, \mathcal{G}) = \Gamma(\mathcal{C}, \mathcal{H})
\quad\text{and}\quad
\Hom_\mathcal{O}(\mathcal{F}, \mathcal{G}_\lambda) =
\Gamma(\mathcal{C}, \mathcal{H}_\lambda)
$$
by construction. By Lemma \ref{lemma-hom-finite-presentation-colimit} we have
$\mathcal{H} = \colim \mathcal{H}_\lambda$. Thus the lemma follows from
Sites, Lemma \ref{sites-lemma-directed-colimits-global-sections}.
\end{proof}









